% subsection: 連立型

\subsection{連立型}\label{節：連立型}
  \begin{bookpoint}
    \textbf{和と差を追う}\par
    和・差など,一つずつ更新できる組合せを選ぶ.元の二つの\bookterm{sequence}{数列}へ戻せる関係と初期値も確かめる.
  \end{bookpoint}
  \subsubsection{和と差で分ける}
    二つの式の係数が入れ替わっているときは,まず足し引きを試そう.
    和と差をそれぞれ求め,元の二つの数列へ戻すところまで計算する.

    \noindent \bookblocklink{example-system-1}{problem}{answer}{【例題】}\par\nobreak
    \begin{enumerate}[label=(\arabic*),parsep=3pt,format=\bookitemlink{example-system-1}{problem}{answer}]
      \item $\displaystyle a_1=1,\,b_1=1.\ \begin{cases}a_{n+1}=3a_n+b_n\\ b_{n+1}=a_n+3b_n\end{cases}$
    \end{enumerate}

    \noindent \bookblocklink{example-system-1}{hint}{problem}{【方針】}\par\nobreak
    \begin{enumerate}[label=(\arabic*),parsep=3pt,format=\bookitemlink{example-system-1}{hint}{problem}]
      \item 係数$3,1$が入れ替わるので,和と差がそれぞれ等比型になります.和と差の初項を求めてから,$a_n,b_n$へ戻します.
    \end{enumerate}

    \noindent \bookblocklink{example-system-1}{answer}{problem}{【解答】}\par\nobreak
    \begin{enumerate}[label=(\arabic*),parsep=3pt,format=\bookitemlink{example-system-1}{answer}{problem}]
      \item $\displaystyle a_n=4^{n-1},\qquad b_n=a_n=4^{n-1}$
    \end{enumerate}

    \noindent \bookblocklink{example-system-1}{solution}{problem}{【解法】}\par\nobreak
    \begin{enumerate}[label=(\arabic*),parsep=3pt,format=\bookitemlink{example-system-1}{solution}{problem}]
      \item 2つの式を辺々加えると,
                  \begin{align*}
                    a_{n+1}+b_{n+1}
                      &=(3a_n+b_n)\\
                      &\quad +(a_n+3b_n)\\
                      &=4(a_n+b_n)
                  \end{align*}
                  よって,
                  \[
                    a_n+b_n=(a_1+b_1)4^{n-1}=2\cdot4^{n-1}. \booktag{sim-example-sum}{1}
                  \]

                  また,2つの式を辺々引くと,
                  \begin{align*}
                    a_{n+1}-b_{n+1}
                      &=(3a_n+b_n)\\
                      &\quad -(a_n+3b_n)\\
                      &=2(a_n-b_n)
                  \end{align*}
                  よって,
                  \[
                    a_n-b_n=(a_1-b_1)2^{n-1}=0. \booktag{sim-example-difference}{2}
                  \]

                  \(\bookeqbackref{sim-example-sum}{(1)}+\bookeqbackref{sim-example-difference}{(2)}\)をすると,
                  \[
                    2a_n=2\cdot4^{n-1}
                  \]
                  なので,
                  \[
                    a_n=4^{n-1}
                  \]

                  また,\(\bookeqref{sim-example-sum}{(1)}-\bookeqref{sim-example-difference}{(2)}\)をすると,
                  \[
                    2b_n=2\cdot4^{n-1}
                  \]
                  なので,
                  \[
                    b_n=4^{n-1}
                  \]
    \end{enumerate}

    \bookstep{係数が対称な場合へ広げる}
    一般形が
    \[
      \begin{cases}
        a_{n+1}=pa_{n}+qb_{n} \\
        b_{n+1}=qa_{n}+pb_{n}
      \end{cases}
    \]
    の形で表される\bookterm{recurrence}{漸化式}の解き方について考えよう.
    二つの式では係数$p,q$の位置が入れ替わっている.
    足すと両方の係数が$p+q$になり,引くと$p-q$とその符号反転が現れる.
    そこで,和と差がそれぞれ単独で更新されないか調べる.
    実際にすすめていくと,次の二つの等比型になる.
    以下の冪の式は$p+q$と$p-q$が零でない場合に使う.零になる方の和または差は,$n\ge2$で$0$となるので,初項と分けて扱う.
    \begin{align*}
      &a_{n+1}+b_{n+1}=(p+q)(a_{n}+b_{n})\\
      &a_n+b_n=(a_1+b_1)(p+q)^{n-1}\booktag{sim-general-sum}{1}
    \end{align*}
    また,
    \begin{align*}
      &a_{n+1}-b_{n+1}=(p-q)(a_{n}-b_{n})\\
      &a_n-b_n=(a_1-b_1)(p-q)^{n-1}\booktag{sim-general-difference}{2}
    \end{align*}
    $\bookeqbackref{sim-general-sum}{(1)}+\bookeqbackref{sim-general-difference}{(2)}$をすると
\BookMathLayout{
\[
\begin{aligned}
2a_n&=(a_1+b_1)(p+q)^{n-1}\\
&\quad+(a_1-b_1)(p-q)^{n-1}\\
a_n&=\frac{1}{2}\bigl\{(a_1+b_1)(p+q)^{n-1}\\
&\qquad+(a_1-b_1)(p-q)^{n-1}\bigr\}
\end{aligned}
\]
}{
\[
\begin{aligned}
2a_n&=(a_1+b_1)(p+q)^{n-1}\\
&\quad+(a_1-b_1)(p-q)^{n-1}\\
a_n&=\frac{1}{2}\bigl\{(a_1+b_1)(p+q)^{n-1}\\
&\qquad+(a_1-b_1)(p-q)^{n-1}\bigr\}
\end{aligned}
\]
}
    $\bookeqref{sim-general-sum}{(1)}-\bookeqref{sim-general-difference}{(2)}$をすると
\BookMathLayout{
\[
\begin{aligned}
2b_n&=(a_1+b_1)(p+q)^{n-1}\\
&\quad-(a_1-b_1)(p-q)^{n-1}\\
b_n&=\frac{1}{2}\bigl\{(a_1+b_1)(p+q)^{n-1}\\
&\qquad-(a_1-b_1)(p-q)^{n-1}\bigr\}
\end{aligned}
\]
}{
\[
\begin{aligned}
2b_n&=(a_1+b_1)(p+q)^{n-1}\\
&\quad-(a_1-b_1)(p-q)^{n-1}\\
b_n&=\frac{1}{2}\bigl\{(a_1+b_1)(p+q)^{n-1}\\
&\qquad-(a_1-b_1)(p-q)^{n-1}\bigr\}
\end{aligned}
\]
}
    和と差を追う場合が解けた.次は,係数が対称でない場合へ広げよう.

  \subsubsection{重みを付けて分ける}
    連立一般型の\bookterm{recurrence}{漸化式}は,
    \[
      \begin{cases}
        a_{n+1}=pa_{n}+qb_{n} \\
        b_{n+1}=ra_{n}+sb_{n}
      \end{cases}
    \]
    で表される\bookterm{recurrence}{漸化式}である.
    今回も\bookterm{geometric}{等比数列}になる量を探そう.
    $a_n$と$b_n$を直接追う代わりに,二つを合わせた量を追う.
    単に二つの式を足すと,
\BookMathLayout{\[
      \begin{array}{rrl}
        & a_{n+1}={} &pa_{n}+qb_{n}\\
        +\,) & b_{n+1}={} & ra_{n}+sb_{n}\\ \hline
        & a_{n+1}+b_{n+1}={} & (p+r)a_{n}+(q+s)b_{n}
      \end{array}
    \]}{
\[
\setlength{\arraycolsep}{2pt}
\begin{array}{rrl}
&a_{n+1}={}&pa_n+qb_n\\
+\,)&b_{n+1}={}&ra_n+sb_n\\\hline
&a_{n+1}+b_{n+1}={}&(p+r)a_n+(q+s)b_n
\end{array}\]
}
    となる.$p+r=q+s$なら和が等比型になるが,一般には一致しない.
    そこで,$a_n,b_n$に重み$\alpha,\beta$を付けて,
    \[
      \alpha a_{n+1}+\beta b_{n+1}
      =\gamma(\alpha a_n+\beta b_n)
    \]
    を満たすようにしてやればよい.
    実際に左辺を計算すると,
\BookMathLayout{\begin{align*}
      \alpha a_{n+1}+\beta b_{n+1}&=\alpha(pa_n+qb_n)+\beta(ra_n+sb_n)\\
      &=(p\alpha+r\beta)a_n+(q\alpha+s\beta)b_n
    \end{align*}}{
\[
\begin{aligned}
\alpha a_{n+1}+\beta b_{n+1}
&=\alpha(pa_n+qb_n)\\ &\quad+\beta(ra_n+sb_n)\\
&=(p\alpha+r\beta)a_n\\ &\quad+(q\alpha+s\beta)b_n
\end{aligned}
\]
}
    一方,
    \[
      \gamma(\alpha a_n+\beta b_n)=\gamma\alpha a_n+\gamma\beta b_n
    \]
    である.
    したがって,
    \[
    \begin{cases}
      p\alpha+r\beta=\gamma\alpha\\
      q\alpha+s\beta=\gamma\beta
    \end{cases}
    \]
    となるように$\alpha,\beta,\gamma$を決めればよい.

    ここで,いきなり3つの文字をすべて求めようとすると少し複雑に見える.
    そこで,まず$\gamma$について考える.

    上の2式を変形すると,
    \[
    \begin{cases}
      (p-\gamma)\alpha+r\beta=0\\
      q\alpha+(s-\gamma)\beta=0
    \end{cases}
    \]
    である.
    $\alpha=\beta=0$では何も取り出せないので,少なくとも一方が零でない組を探す.
    第1式を$s-\gamma$倍し,第2式を$r$倍して引くと,$\beta$が消えて
    \[
      \{(p-\gamma)(s-\gamma)-qr\}\alpha=0
    \]
    となる.同様に,第2式を$p-\gamma$倍して第1式の$q$倍を引けば,
    \[
      \{(p-\gamma)(s-\gamma)-qr\}\beta=0
    \]
    となる.$\alpha,\beta$は同時に零ではないので,
    \[
      (p-\gamma)(s-\gamma)-qr=0
    \]
    でなければならない.
    整理すると,
    \[
      \gamma^2-(p+s)\gamma+(ps-qr)=0
    \]
    となる.
    この二次方程式を,この連立\bookterm{recurrence}{漸化式}の\bookterm{characteristic}{特性方程式}と呼ぶ.
    \bookterm{characteristic}{特性方程式}の解を$\gamma_1,\gamma_2$とする.
    まず$\gamma=\gamma_1$を上の連立方程式に代入し,
    \[
      \begin{cases}
        (p-\gamma_1)\alpha+r\beta=0\\
        q\alpha+(s-\gamma_1)\beta=0
      \end{cases}
    \]
    を満たす$\alpha,\beta$を1組求める.
    すると,
    \[
      c_n=\alpha a_n+\beta b_n
    \]
    と置くことで,
    \begin{align*}
      c_{n+1}&=\alpha a_{n+1}+\beta b_{n+1}\\
      &=\gamma_1(\alpha a_n+\beta b_n)\\
      &=\gamma_1c_n
    \end{align*}
    となる.
    したがって,
    \[
      c_n=c_1\gamma_1^{\,n-1}
    \]
    という\bookterm{geometric}{等比数列}が得られる.ただし$\gamma_1=0$なら,$c_1$を別に保ち,$n\ge2$では$c_n=0$とする.
    同様に,\,$\gamma=\gamma_2$に対しても$\alpha,\beta$を求めれば,
    \[
      d_n=\alpha' a_n+\beta' b_n
    \]
    について
    \[
      d_n=d_1\gamma_2^{\,n-1}
    \]
    という\bookterm{geometric}{等比数列}が得られる.$\gamma_2=0$の場合も同じく初項と分ける.
    なお,ここで得られた2つの\bookterm{geometric}{等比数列}から$a_n,b_n$を求めるには,
    \[
      \begin{cases}
      c_n=\alpha a_n+\beta b_n\\
      d_n=\alpha'a_n+\beta'b_n
      \end{cases}
    \]
    を$a_n,b_n$について解けばよい.
    ただし,この連立方程式から$a_n,b_n$を一意に求めるには,係数の組$(\alpha,\beta)$と$(\alpha',\beta')$が比例しないこと,すなわち$\alpha\beta'-\alpha'\beta\ne0$が必要である.
    異なる二つの特性根なら,この条件を満たす.
    もし二組の重みが比例すれば,同じ非零の重み$(\alpha,\beta)$が両方の係数条件を満たす.
    その条件を引くと$(\gamma_1-\gamma_2)\alpha=(\gamma_1-\gamma_2)\beta=0$となる.
    $\gamma_1\ne\gamma_2$なら$\alpha=\beta=0$になり,非零の重みを選んだことに矛盾する.
    ただし,異なる根は十分条件であり,必要条件ではない.
    重解では,比例しない二組を取れるかどうかで手順が分かれる.\footnote{例えば$a_{n+1}=2a_n,\,b_{n+1}=2b_n$では$(1,0),(0,1)$を選べる.
    一方,比例しない二組を取れない場合は,一方の量を先に求め,残る式へ代入する.\bookproblemref{practice-basic}{25}{標準問25}では,その式を指数の付加項として解く.\booknavref{節：二次特性方程式型}{隣接3項型の重解}とも比べられる.}
    このように,\,$a_n,b_n$をそれぞれ直接求めるのではなく,\,$a_n,b_n$の和の形を作ることで,連立していた\bookterm{recurrence}{漸化式}を\bookterm{geometric}{等比数列}へと分離することができる.\footnote{また,$\,\gamma_1,\gamma_2$が虚数になる場合もこの手法は形式的には成り立つが,実数の形に書き直すには\sectionref{節：虚数解型隣接3項間漸化式}で扱う手法が必要になる.}


    \noindent \hyperlink{book:example-system-1:answer:2}{\textbf{【例題】}}\par\nobreak
    \begin{enumerate}[start=2,label=(\arabic*),parsep=3pt,format=\bookitemlink{example-system-1}{problem}{answer}]
      \item $\displaystyle a_1=1,\,b_1=0,\ \begin{cases}a_{n+1}=3a_n+b_n\\ b_{n+1}=2a_n+2b_n\end{cases}$
    \end{enumerate}

    \noindent \hyperlink{book:example-system-1:problem:2}{\textbf{【方針】}}\par\nobreak
    \begin{enumerate}[start=2,label=(\arabic*),parsep=3pt,format=\bookitemlink{example-system-1}{hint}{problem}]
      \item $\gamma$を求めてから重みを選ぶと,$a_n-b_n$と$2a_n+b_n$をそれぞれ単独で更新できます.
    \end{enumerate}

    \noindent \hyperlink{book:example-system-1:problem:2}{\textbf{【解答】}}\par\nobreak
    \begin{enumerate}[start=2,label=(\arabic*),parsep=3pt,format=\bookitemlink{example-system-1}{answer}{problem}]
      \item $\displaystyle a_n=\frac{2\cdot4^{n-1}+1}{3},\qquad b_n=\frac{2\cdot4^{n-1}-2}{3}$
    \end{enumerate}

    \noindent \hyperlink{book:example-system-1:problem:2}{\textbf{【解法】}}\par\nobreak
    \begin{enumerate}[start=2,label=(\arabic*),parsep=3pt,format=\bookitemlink{example-system-1}{solution}{problem}]
      \item $\displaystyle a_1=1,\,b_1=0,\ \begin{cases}a_{n+1}=3a_n+b_n\\ b_{n+1}=2a_n+2b_n\end{cases}$\\
          $\alpha a_{n+1}+\beta b_{n+1}$を計算すると,
          \begin{align*}
            \alpha a_{n+1}+\beta b_{n+1}
              &=\alpha(3a_n+b_n)\\
              &\quad +\beta(2a_n+2b_n)\\
              &=(3\alpha+2\beta)a_n\\
              &\quad +(\alpha+2\beta)b_n
          \end{align*}
          これが
          \[
            \gamma(\alpha a_n+\beta b_n)
            =\gamma\alpha a_n+\gamma\beta b_n
          \]
          と一致するためには,
          \[
          \begin{cases}
            (3-\gamma)\alpha+2\beta=0\\
            \alpha+(2-\gamma)\beta=0
          \end{cases}
          \]
          でなければならない.

          $\alpha,\beta$が同時に$0$ではないため,
          \[
            (3-\gamma)(2-\gamma)-2=0
          \]
          より,
          \begin{align*}
            &\gamma^2-5\gamma+4=0\\
            &(\gamma-1)(\gamma-4)=0
          \end{align*}
          したがって,
          \[
            \gamma_1=1,\qquad\gamma_2=4
          \]
          である.

          まず$\gamma=1$とすると,
          \[
          \begin{cases}
            2\alpha+2\beta=0\\
            \alpha+\beta=0
          \end{cases}
          \]
          となるので,
          \[
            \alpha=1,\qquad\beta=-1
          \]
          と取ることができる.

          そこで,
          \[
            c_n=a_n-b_n
          \]
          とおくと,
          \begin{align*}
            c_{n+1}
              &=a_{n+1}-b_{n+1}\\
              &=(3a_n+b_n)-(2a_n+2b_n)\\
              &=a_n-b_n\\
              &=c_n
          \end{align*}
          また,
          \[
            c_1=a_1-b_1=1
          \]
          であるから,
          \[
            c_n=1
          \]
          したがって,
          \[
            a_n-b_n=1
          \]

          次に$\gamma=4$とすると,
          \[
          \begin{cases}
            -\alpha+2\beta=0\\
            \alpha-2\beta=0
          \end{cases}
          \]
          となるので,
          \[
            \alpha=2,\qquad\beta=1
          \]
          と取ることができる.

          そこで,
          \[
            d_n=2a_n+b_n
          \]
          とおくと,
          \begin{align*}
            d_{n+1}
              &=2a_{n+1}+b_{n+1}\\
              &=2(3a_n+b_n)+(2a_n+2b_n)\\
              &=8a_n+4b_n\\
              &=4(2a_n+b_n)\\
              &=4d_n
          \end{align*}
          また,
          \[
            d_1=2a_1+b_1=2
          \]
          であるから,
          \[
            d_n=2\cdot4^{n-1}
          \]

          以上より,
          \[
          \begin{cases}
            a_n-b_n=1\\
            2a_n+b_n=2\cdot4^{n-1}
          \end{cases}
          \]
          となる.

          2式を辺々加えると,
          \[
            3a_n=1+2\cdot4^{n-1}
          \]
          であるから,
          \[
            a_n=\frac{2\cdot4^{n-1}+1}{3}
          \]
          また,
          \[
            b_n=a_n-1
          \]
          より,
          \[
            b_n=\frac{2\cdot4^{n-1}-2}{3}
          \]
    \end{enumerate}

    \bookstep{和と差は重みの一例}
    \bookproblemref{example-system-1}{1}{最初の例題}へ戻ろう.
    和と差を思い付かなかったとしても,今の係数合わせから同じ量を選べる.
    $p=s=3,\ q=r=1$を特性方程式へ入れると,
    \[
      \gamma^2-6\gamma+8=(\gamma-2)(\gamma-4)=0
    \]
    となる.$\gamma=2$では重みの条件は二式とも$\alpha+\beta=0$となるので,
    $(\alpha,\beta)=(1,-1)$を選べる.
    $\gamma=4$では二式とも$\alpha=\beta$となり,$(1,1)$を選べる.
    従って取り出す量は,
    \[
      c_n=a_n-b_n,\qquad d_n=a_n+b_n
    \]
    であり,更新は$c_{n+1}=2c_n,\ d_{n+1}=4d_n$となる.
    初項も$c_1=0,\ d_1=2$なので,最初の解法と同じ二つの式が得られる.
    対称な係数を見て足し引きする方法は短い.
    重みを条件から求める方法は,係数が対称でない例題(2)にも使える.
    二つの見方は,別の答えを作るのではなく,同じ単独の更新を取り出している.
